% Reduced extract from Jiri Lebl, Tasty Bits of Several Complex Variables.
% License: CC BY-SA 4.0. See SOURCES.md for attribution and source revision.
\documentclass{book}
\usepackage{amsmath,amssymb}
\newcommand{\C}{\mathbb{C}}
\newcommand{\R}{\mathbb{R}}
\newcommand{\N}{\mathbb{N}}
\newcommand{\abs}[1]{\lvert #1 \rvert}

\begin{document}
\chapter{Holomorphic Functions in Several Variables}
\section{Motivation, single variable, and Cauchy's formula}\label{sec:motivation}

We start with some standard notation. We use $\C$ for complex numbers,
$\R$ for real numbers, and $\N = \{1,2,3,\ldots\}$ for natural numbers.
As complex analysis deals with complex numbers, perhaps we should begin with
$\sqrt{-1}$. Start with the real numbers, $\R$, and add $\sqrt{-1}$ into our
field. Call this square root $i$, and write the complex numbers, $\C$, by
identifying $\C$ with $\R^2$ using
\begin{equation*}
z = x + iy,
\end{equation*}
where $z \in \C$ and $(x,y) \in \R^2$.

Given a complex number $z$, its "opposite" is the \emph{complex conjugate}
of $z$ and is defined as
\begin{equation*}
\bar{z} \overset{\text{def}}{=} x - iy.
\end{equation*}
The size of $z$ is measured by the \emph{modulus}:
\begin{equation*}
\abs{z} \overset{\text{def}}{=} \sqrt{z \bar{z}} = \sqrt{x^2 + y^2}.
\end{equation*}

A nonzero polynomial in $z$ is an expression
\begin{equation*}
P(z) = \sum_{k=0}^d c_k\,z^k,
\end{equation*}
where $c_k \in \C$ and $c_d \neq 0$. The number $d$ is called the
\emph{degree} of the polynomial $P$.
\end{document}
